% Exact per-target-lead identifiability (exact pseudo-inverse) vs the deployed rho-truncated
% numerical map. No residual-independence claim on real data; the strict lower bound is stated
% only for the synthetic model (Proposition 2).

\newtheorem{theoremc}{Proposition}
\newtheorem{propc}[theoremc]{Proposition}

% ---------------------------------------------------------------- setup
\noindent\textbf{Model.} Within segment $s$ write the population 12-lead vector as
$\bm L_s = \bm\mu_s + \bm M_s\bm d + \bm r_s$, where $\bm M_s\in\mathbb R^{12\times3}$
is the (population-estimated, PCA) dipolar basis, $\bm d$ the dipole coordinates, and
$\bm r_s$ the non-dipolar residual. Observing $S$ gives $\bm y_S=\bm{\mathrm{Sel}}_S\bm L_s
+\bm n$, with $\bm M_{s,S}=\bm{\mathrm{Sel}}_S\bm M_s$. Write the exact Moore--Penrose
pseudo-inverse $\bm M_{s,S}^{+}$ and the exact row-space projector
$\bm P_{\mathrm{obs}}=\bm M_{s,S}^{+}\bm M_{s,S}$. For numerical deployment we also use the
\emph{rank-$r$ truncation} at relative tolerance $\varrho$,
$\bm M_{s,S}^{(\varrho)}=\bm U_r\bm\Sigma_r\bm V_r^{\top}$ keeping singular triplets with
$\sigma_i>\varrho\,\sigma_{\max}$, with pseudo-inverse $\bm M_{s,S}^{(\varrho)+}$ and projector
$\bm P_{\mathrm{obs}}^{(\varrho)}=\bm V_r\bm V_r^{\top}$.

% ---------------------------------------------------------------- exact theorem
\begin{theoremc}[Exact per-target-lead dipolar identifiability and conditioning]
\label{thm:perlead}
For each target lead $\ell$ define
$\eta_{s,\ell}(S)=\big\|\bm e_\ell^{\top}\bm M_s(\bm I-\bm P_{\mathrm{obs}})\big\|_2$ and
$\kappa_{s,\ell}(S)=\big\|\bm e_\ell^{\top}\bm M_s\bm M_{s,S}^{+}\big\|_2$ from the \emph{exact}
$\bm M_{s,S}$.
\begin{enumerate}
\item[(i)] \emph{Identifiability.} $\eta_{s,\ell}(S)=0$ iff $\bm e_\ell^{\top}\bm M_s$ lies in
the exact row space of $\bm M_{s,S}$; then $\bm e_\ell^{\top}\bm M_s\bm d$ is a deterministic
function of $\bm M_{s,S}\bm d$---two dipoles with the same observation have the same lead-$\ell$
dipolar value, at any SNR. If $\eta_{s,\ell}(S)>0$ there is a $\bm\delta\in\ker\bm M_{s,S}$
(an unobserved dipole direction) with $\bm e_\ell^{\top}\bm M_s\bm\delta\neq0$; then $\bm d$ and
$\bm d+t\bm\delta$ produce identical observations for all $t$ while lead $\ell$'s dipolar value
changes, so it is not identifiable at any SNR.
\item[(ii)] \emph{Conditioning.} On the identifiable part, the lead-$\ell$ error from
observation noise and observed non-dipolar residual obeys
$\big|\bm e_\ell^{\top}\bm M_s\bm M_{s,S}^{+}(\bm{\mathrm{Sel}}_S\bm r_s+\bm n)\big|
\le \kappa_{s,\ell}(S)\,\big(\|\bm{\mathrm{Sel}}_S\bm r_s\|_2+\|\bm n\|_2\big).$
\end{enumerate}
The global $\kappa_s(S)=\|\bm M_s\bm M_{s,S}^{+}\|_2$ (spectral norm) upper-bounds every per-lead
$\kappa_{s,\ell}(S)$ (a row norm never exceeds the spectral norm; it is not their maximum).
\end{theoremc}

\noindent\emph{Exact vs.\ $\varrho$-identifiability (numerical deployment).} The deployed map
computes $\eta^{(\varrho)}_{s,\ell},\kappa^{(\varrho)}_{s,\ell}$ from the truncation
$\bm P_{\mathrm{obs}}^{(\varrho)}$, so a singular direction observed only through a value below
$\varrho\,\sigma_{\max}$ is treated as unobserved. Then $\eta^{(\varrho)}_{s,\ell}>0$ means lead
$\ell$ is \emph{not stably recoverable at resolution $\varrho$} (outside the retained numerical
row space); it does \emph{not} by itself imply the exact matrix has a null direction---a
tiny-but-nonzero singular value is exactly identifiable ($\eta=0$) yet $\varrho$-unidentifiable
($\eta^{(\varrho)}>0$), as we verify numerically (\texttt{tests/test\_exact\_vs\_rho.py}). We therefore call the
$\eta^{(\varrho)}$ heatmap the \emph{$\varrho$-identifiability} map; for well-separated
configurations the exact and $\varrho$ verdicts coincide, and near-rank-deficient $S$ are
reported with a $\varrho$-sweep and record-bootstrap CIs. \textbf{The limb-6 flagship is exact:}
the six limb leads are exact linear combinations of I and II (Einthoven/Goldberger), so
$\bm M_{s,S}$ has exact rank~2 and precordial ST is \emph{exactly} non-identifiable from limb
leads at any SNR, independent of $\varrho$.

\noindent\emph{Relation to classical estimability.} Part~(i) is the classical estimability
criterion (Gauss--Markov / \cite{scheffe1959anova}) for the linear functional
$\bm e_\ell^{\top}\bm M_s\bm d$; we claim no novelty for this identity. The contribution
(Sec.~\ref{sec:exp}) is its use as a graded, empirically-estimated recoverability map with
per-record calibrated error rates (Sec.~\ref{sec:calib}).

% ---------------------------------------------------------------- residual (honest)
\noindent\textbf{The non-dipolar residual (no independence claim on real data).}
Let $\bm m_s(\bm y_S)=\mathbb E[\bm r_s\mid\bm y_S]$ and
$\bm u_s=\bm r_s-\bm m_s(\bm y_S)$, so $\mathbb E[\bm u_s\mid\bm y_S]=\bm 0$. We do
\emph{not} assume $\bm u_s\perp\bm y_S$. Under squared loss the Bayes-optimal reconstruction
error of the (vector) residual is
\[
\inf_{\hat{\bm r}}\ \mathbb E\big\|\bm r_s-\hat{\bm r}(\bm y_S)\big\|_2^2
=\mathbb E\big[\operatorname{tr}\operatorname{Cov}(\bm r_s\mid\bm y_S)\big].
\]
Part of $\bm r_s$ (the \emph{empirically predictable residual}) is recovered by a
predictor trained on $\bm y_S$ and metered with distribution-free calibrated intervals;
the remainder is an \emph{achievability gap}, established by a held-out achievability
analysis for a stated predictor family---\emph{not} declared unrecoverable a priori.

\begin{propc}[Strict independence limit -- synthetic model only]
\label{prop:synth}
If the data are generated so that a component $\bm u$ is statistically independent of
$\bm y_S$ (an explicit synthetic construction), then $p(\bm u\mid\bm y_S)=p(\bm u)$, the
Bayes estimate of $\bm u$ under squared loss is its prior mean, and
$\mathbb E\|\widehat{\bm u}-\bm u\|_2^2\ge\operatorname{tr}\operatorname{Cov}(\bm u)$ for any
estimator. This strict limit is invoked \emph{only} for the synthetic model; on real ECG we
make no such independence claim.
\end{propc}
